\section{Proposed Model(s) and Baseline}
% \begin{itemize}
%     \item Which simple baseline model will you compare your solution to?
%     \item Which ``proper'' model(s) will you use to get a well-performing ML system?
%     \item Why are you choosing these models?
%     \item How will you train your model(s)?
%     \item What hyperparameter tuning strategy will you implement?
%     \item Which hyperparameters will you tune?
%     \item If you use an SVM, what kernel function will you choose and why?
%     \item \textbf{Required:} At least 1 baseline and at least 1 ``full'' model.
% \end{itemize}

We plan to:
\begin{itemize}
    \item \textbf{Baseline:} FFT-based features with an SVM, or a similarly simple classical model.
    \item \textbf{Deep learning model:} still to be decided, but it will likely be trained on raw or lightly processed windows.
    \item \textbf{Loss:} triplet loss.
    \item \textbf{Model comparison:} accuracy, precision, recall, F1 score, confusion matrix, and $K$-fold comparison.
    \item \textbf{Hyperparameter search:} grid search.
    \item \textbf{Regularization techniques:} dropout, L2 regularization, early stopping based on F1, batch normalization, layer normalization, and Adam for optimization.
    \item \textbf{Hyperparameters to tune:} dropout, L2 strength, layer count, and node count.
    \item \textbf{Deployment goal:} compress the model to run on a phone through pruning, compression, quantization, or related methods.
    \item \textbf{Application goal:} build a web or phone app that records gyroscopic data live, feeds it to the model, and returns live predictions.
\end{itemize}

